\origpage{563}{N. Robertson, P. Seymour / Journal of Combinatorial Theory, Series B 102 (2012) 530--563}
\noindent $d(v^*, D) \ge \theta > \theta' + 2$. Also, by (10), $d(v, D) \ge \theta - \theta' - 2 > h$. By hypothesis (xiii), there exists $A' \in \mathcal{A}$ such that $d(v, v(A')) \le h$ and $(A', \pi(A') + \omega)$ has the same $\delta$-folio as $(A, \pi(A) + \omega)$. Hence, by (6), $d(v, D(A')) \le h + 1 = \eta$. Finally,
\begin{equation*}
\theta - \theta' - 2 \le d(v, D) \le d(v, v(A')) + d(v(A'), D) \le h + d(v(A'), V(C_0))
\end{equation*}
and so $d(v(A'), V(C_0)) \ge \theta - \theta' - 2 - h \ge 5$. Hence $A' \in \mathcal{A}'$. This proves (13).

Consequently, all the hypotheses of 8.2 hold with $\mathcal{A}$ and $\theta$ replaced by $\mathcal{A}'$ and $\theta'$, and it follows from 8.2 that there is a model of $(H, \chi)$ in $(G \setminus \{v^*\}, \omega)$. We deduce that the $\delta$-folio of $(G, \omega)$ is a subset of the $\delta$-folio of $(G \setminus \{v^*\}, \omega)$, and we therefore have equality, since the reverse inclusion is trivial. The result follows. $\square$

Finally, a few words on deriving Theorem 10.2 of [5] from 8.3. In the language of [5] we have a wall with an $h$-homogeneous subwall of height $\theta$ ($\theta$ replaces the $f(h)$ of [5]). Take $\Gamma$ to be the original wall, and let $C_0$ be its perimeter. This wall has height at least $\theta$ since it has a subwall of height $\theta$, and hence it contains the $\theta \times \theta$ grid as a minor, and from theorems 6.1 and 7.3 of [3] it therefore has a tangle $\mathcal{T}$ of order $\theta$. Let $\Pi$ be the set of corners of $\Gamma$, and let $\mathcal{A}$ be the set of graphs called $\tilde{A}$ in the final section of [5]. Let $\pi(A) + \omega$ be the ``attachment sequence'' of $A$ in the language of [5]. Let $D$ be the closed disc with boundary the perimeter of the $h$-homogeneous subwall including the ``infinite'' region of $\Gamma$. Then hypotheses (i)--(xiv) all are satisfied (for (xiii), we use that the subwall is $h$-homogeneous). Consequently, Theorem 10.2 of [5] is true.

\sectionhead{References}

[1] N. Robertson, P.D. Seymour, Graph minors. VI. Disjoint paths across a disc, J. Combin. Theory Ser. B 41 (1986) 115--138. [2] N. Robertson, P.D. Seymour, Graph minors. VII. Disjoint paths on a surface, J. Combin. Theory Ser. B 45 (1988) 212--254. [3] N. Robertson, P.D. Seymour, Graph minors. X. Obstructions to tree-decomposition, J. Combin. Theory Ser. B 52 (1991) 153--190. [4] N. Robertson, P.D. Seymour, Graph minors. XII. Distance on a surface, J. Combin. Theory Ser. B 64 (1995) 240--272. [5] N. Robertson, P.D. Seymour, Graph minors. XIII. The disjoint paths problem, J. Combin. Theory Ser. B 63 (1995) 65--110. [6] N. Robertson, P.D. Seymour, Graph minors. XIV. Extending an embedding, J. Combin. Theory Ser. B 65 (1995) 23--50. [7] N. Robertson, P.D. Seymour, Graph minors. XXI. Graphs with unique linkages, J. Combin. Theory Ser. B 99 (2009) 583--616. [8] P.D. Seymour, R. Thomas, Graph searching, and a min-max theorem for tree-width, J. Combin. Theory Ser. B 58 (1993) 22--33.
